\section{API Design und Controller}
Die Kommunikation mit dem Frontend erfolgt \"uber eine RESTful API. Die acht Controller sind thematisch gruppiert und folgen REST-Konventionen (\texttt{GET}, \texttt{POST}, \texttt{PUT}, \texttt{PATCH}, \texttt{DELETE}). Nahezu alle Endpunkte sind hinter dem \texttt{[Authorize]}-Attribut gesch\"utzt, d.h. sie erfordern eine aktive Cookie-Session. Ausnahmen sind die \"offentlichen Authentifizierungs-Endpunkte (Registrierung, Login, E-Mail-Best\"atigung).

\subsection{Ausgew\"ahlte API-Endpunkte}

\begin{table}[H]
\centering
\begin{tabular}{|l|p{5.5cm}|l|p{4cm}|}
\hline
\textbf{Methode} & \textbf{Route} & \textbf{Auth} & \textbf{Beschreibung} \\
\hline
\multicolumn{4}{|l|}{\textbf{api/users}} \\
\hline
\texttt{POST}   & \texttt{/}              & Anonym & Registrierung \\
\texttt{POST}   & \texttt{/login}         & Anonym & Login, setzt HttpOnly-Cookie \\
\texttt{POST}   & \texttt{/logout}        & Auth   & Logout, l\"oscht Session \\
\texttt{GET}    & \texttt{/me}            & Auth   & Eigenes Profil abrufen \\
\texttt{DELETE} & \texttt{/\{id\}}        & Admin  & Nutzer l\"oschen \\
\hline
\multicolumn{4}{|l|}{\textbf{api/documents}} \\
\hline
\texttt{GET}    & \texttt{/}              & Auth   & Eigene Dokumente (paginiert) \\
\texttt{POST}   & \texttt{/}              & Auth   & Dokument hochladen (startet Indexierungs-Job) \\
\texttt{DELETE} & \texttt{/\{id\}}        & Auth   & Dokument inkl. MinIO-Datei l\"oschen \\
\texttt{GET}    & \texttt{/download/\{id\}} & Auth & Datei-Stream herunterladen \\
\texttt{POST}   & \texttt{/\{id\}/share}  & Auth   & Dokument per E-Mail teilen \\
\texttt{GET}    & \texttt{/search}        & Auth   & Volltext- und Tag-Suche \\
\hline
\multicolumn{4}{|l|}{\textbf{api/chat}} \\
\hline
\texttt{POST}   & \texttt{/message}       & Auth   & Nachricht senden, l\"ost RAG-Pipeline aus \\
\texttt{GET}    & \texttt{/sessions}      & Auth   & Chat-Sitzungen abrufen \\
\texttt{POST}   & \texttt{/sessions}      & Auth   & Neue Chat-Sitzung erstellen \\
\hline
\multicolumn{4}{|l|}{\textbf{api/flashcards \& api/quizzes}} \\
\hline
\texttt{POST}   & \texttt{/flashcards/generate} & Auth & Lernkarten per KI generieren \\
\texttt{POST}   & \texttt{/quizzes/generate}    & Auth & Quiz per KI generieren \\
\hline
\multicolumn{4}{|l|}{\textbf{api/tags}} \\
\hline
\texttt{GET}    & \texttt{/}              & Auth   & Alle Tags abrufen \\
\texttt{POST}   & \texttt{/}              & Admin  & Neuen Tag erstellen \\
\hline
\end{tabular}
\caption{Wichtigste API-Endpunkte des EduShelf-Backends}
\label{tab:api_endpoints}
\end{table}
